\documentclass[11pt,a4paper]{article}
\usepackage[utf8]{inputenc}
\usepackage[T1]{fontenc}
\usepackage[brazil]{babel}
\usepackage[margin=2cm]{geometry}
\usepackage{lmodern}
\usepackage{array,tabularx}
\usepackage{colortbl}
\usepackage{booktabs}
\usepackage{xcolor}
\usepackage{tikz}
\usetikzlibrary{arrows.meta}
\usepackage{hyperref}
\hypersetup{colorlinks=true,urlcolor=blue,linkcolor=black}
\setlength{\parindent}{0pt}
\setlength{\parskip}{7pt}
\renewcommand{\arraystretch}{1.3}

\begin{document}

\begin{titlepage}
\centering
\vspace*{0.2cm}
{\fontsize{40}{44}\selectfont\sffamily U\textcolor[HTML]{861D23}{F}\textcolor[HTML]{C9A733}{V}\par}
\vspace{0.2cm}
{\small SERVIÇO PÚBLICO FEDERAL - MINISTÉRIO DA EDUCAÇÃO\par}
\vspace{0.15cm}
{\small UNIVERSIDADE FEDERAL DE VIÇOSA - UFV\par}
\vspace{0.15cm}
{\small CAMPUS FLORESTAL\par}
\vspace{3.5cm}
{\Large\bfseries Trabalho Prático - POO\par}
\vspace{0.4cm}
{\large\bfseries Tartarugas vs. Lixo Marinho\par}
\vspace{0.2cm}
{\normalsize Relatório Técnico - Primeira Entrega Parcial\par}
\vspace{0.4cm}
{\small\textbf{Professor:} Philipe de Freitas Melo\par}
{\small\textbf{Disciplina:} CCF313 -- Programação Orientada a Objetos\par}
\vspace{1.8cm}
\begin{flushright}
Lethícia Gonçalves da Silva - 6571\\[3pt]
Tomás Viana Resende - 6556\\[3pt]
Bruna Lima Ribeiro - 6557\\[3pt]
Helena Santos Perez - 6564\\[3pt]
Otávio Soares Pedroso - 6593\\[3pt]
Pedro Mendes Gonzaga - 6563\\[3pt]
Marina Larissa Xavier Soares - 6581\\[3pt]
Caroline Melo Silva - 6595\\[3pt]
\end{flushright}
\vfill
Florestal - MG\\[5pt]
2026
\end{titlepage}

\clearpage
\section*{1. Proposta do Jogo}
O projeto \textit{Tartarugas vs. Lixo Marinho} constitui um jogo eletrônico do gênero \textit{Tower Defense} bidimensional desenvolvido na linguagem Java para a disciplina CCF313 -- Programação Orientada a Objetos da Universidade Federal de Viçosa (UFV), \textit{Campus} Florestal. A proposta combina conceitos centrais da orientação a objetos com uma temática de conscientização sobre o descarte inadequado de resíduos sólidos em ecossistemas costeiros, ambientando a dinâmica de defesa em uma praia onde espécies de tartarugas marinhas atuam como defensores para impedir o avanço de dejetos antrópicos em direção ao oceano.

\subsection*{Objetivo e Escopo}
O objetivo da partida consiste em defender a costa litorânea ao longo de dez ondas consecutivas de resíduos que percorrem um trajeto pré-fixado. O jogador inicia a partida com 100 conchas (unidade monetária do jogo) e 20 pontos de integridade da praia. A cada dejeto neutralizado pelo ataque das tartarugas, conchas e pontos para o ranking são concedidos ao jogador. Caso a integridade da praia seja reduzida a zero pelo avanço dos resíduos, a partida encerra-se em derrota; a sobrevivência ao término da última onda caracteriza a vitória do jogador.

O escopo do sistema abrange três espécies de tartarugas com comportamentos de combate distintos e três tipos de resíduos com diferentes atributos de deslocamento e resistência. O sistema contempla mecanismos de aquisição e posicionamento de torres, controle de fluxo da partida (início de onda e pausa), e registro de melhores pontuações em memória. Prevê-se a adoção de um mecanismo simplificado de persistência em arquivos de texto em etapa subsequente. O escopo delimita-se a uma aplicação local autônoma, sem necessidade de banco de dados externo ou conexão de rede.

\subsection*{Estratégia e Metodologia de Desenvolvimento}
Adotou-se uma metodologia de desenvolvimento modular e incremental orientada a objetos. Nesta etapa inicial, o foco do trabalho consistiu na modelagem e construção do núcleo de domínio da aplicação, priorizando a hierarquia estrutural de combate (entidades, torres de defesa, projéteis e dejetos marinhos) e o subsistema de pontuação e ordenação de ranking em memória.

Essa estratégia foi definida com o intuito de isolar as regras de negócio centrais e a mecânica balística de camadas periféricas, tais como interfaces de usuário e componentes de persistência física em disco. A separação clara de responsabilidades proporciona baixo acoplamento e alta coesão, viabilizando a validação independente das operações fundamentais por meio de testes unitários automatizados antes da integração com o ciclo de execução contínua (\textit{GameLoop}) e com os módulos de entrada e saída.

\subsection*{Escopo e Status da Primeira Etapa}
A primeira etapa do projeto consolida a infraestrutura conceitual, a modelagem em diagramas de classes e a implementação do núcleo funcional básico do sistema. A Tabela~\ref{tab:status} apresenta a distribuição estruturada entre os componentes já implementados e aqueles projetados para as fases seguintes.

\begin{table}[h!]
\centering
\small
\begin{tabularx}{\textwidth}{|p{0.24\textwidth}|p{0.28\textwidth}|X|}
\hline
\textbf{Subsistema} & \textbf{Componentes} & \textbf{Status e Descrição} \\
\hline
\multicolumn{3}{|c|}{\cellcolor{gray!15}\textbf{Componentes Implementados (Primeira Entrega Parcial)}} \\
\hline
Classe Principal Executável & \texttt{Main.java} & \textbf{Implementada}. Ponto de entrada para demonstração da instanciação de classes e execução do CRUD em console. \\
\hline
Funcionalidade Completa & \texttt{GerenciadorRanking}, \texttt{RegistroRanking} & \textbf{Implementados}. Módulo completo com Create, Read, Update, Delete e ordenação automática (\texttt{Comparable}). \\
\hline
Suíte de Testes Automatizados & \texttt{GerenciadorRankingTest} & \textbf{Implementada}. Conjunto de testes unitários sob o framework JUnit 5 validando o CRUD e exceções. \\
\hline
Hierarquia Base de Combate & \texttt{Entidade}, \texttt{TorreTartaruga}, \texttt{InimigoLixo} & \textbf{Implementadas}. Classes abstratas com encapsulamento de estado ativo, coordenadas e atributos basilares. \\
\hline
Torres Concretas & \texttt{TartarugaDeCouro}, \texttt{TartarugaDePente} & \textbf{Implementadas}. Subclasses com especialização polimórfica de dano, cadência e mecânica de acertos críticos. \\
\hline
Sistema Balístico & \texttt{Projetil} & \textbf{Implementada}. Entidade balística responsável por perseguir alvos e aplicar dano por impacto. \\
\hline
\multicolumn{3}{|c|}{\cellcolor{gray!15}\textbf{Componentes Planejados (Etapas Futuras)}} \\
\hline
Torre Inicial & \texttt{TartarugaVerde} & \textbf{Planejada}. Subclasse de tartaruga com baixo custo para início de partida. \\
\hline
Resíduos Concretos & \texttt{GarrafaPet}, \texttt{RedeFantasma}, \texttt{Microplastico} & \textbf{Planejados}. Subclasses de dejetos marinhos com perfis de velocidade, resistência e conchas. \\
\hline
Tabuleiro e Economia & \texttt{Mapa}, \texttt{CelulaGrid}, \texttt{Jogador}, \texttt{LojaService} & \textbf{Planejados}. Malha de posicionamento, caminho dos resíduos, saldo de conchas e regras de compra. \\
\hline
Motor e Ondas & \texttt{GameLoop}, \texttt{WaveService}, \texttt{CollisionService} & \textbf{Planejados}. Temporização contínua, geração escalonada de resíduos e detecção de colisões. \\
\hline
Interface Gráfica & Interface Gráfica do Usuário (GUI) & \textbf{Planejada}. Telas gráficas simples em JavaFX ou Swing com botões, animações e renderização do mapa. \\
\hline
Persistência Física & \texttt{GerenciadorPersistencia}, \texttt{DadosPartida} & \textbf{Planejados}. Serialização de recordes e salvamento de partidas em arquivos locais de texto plano. \\
\hline
\end{tabularx}
\caption{Quadro comparativo de componentes implementados e planejados.}
\label{tab:status}
\end{table}

\clearpage
\section*{2. Regras e Mecânicas de Jogo}
\subsection*{Mapeamento e Dinâmica de Combate}
O mapa da partida é estruturado sobre uma malha quadriculada bidimensional contendo 12 colunas e 8 linhas. Os resíduos movem-se ao longo de um trajeto fixo e contínuo, iniciando em $(0,3)$, estendendo-se linearmente até $(5,3)$, descendo para $(5,5)$ e prosseguindo horizontalmente até a extremidade $(11,5)$, onde atingem a praia. As células desocupadas que não pertencem ao trajeto estão disponíveis para a instalação de torres de tartaruga pelo jogador.

A classe \texttt{Projetil} é responsável por intermediar os ataques das torres, deslocando-se em direção à posição do dejeto-alvo e aplicando a redução em sua integridade no momento do impacto. Cada dejeto neutralizado pelo combate concede conchas e pontos ao jogador de forma única. Caso um dejeto alcance o término do caminho sem ser destruído, ele é removido da partida, causando o decréscimo de um ponto de vida da praia sem conceder recompensa financeira. Uma onda é dada como concluída quando todos os inimigos previstos forem eliminados ou tiverem chegado ao fim do trajeto.

\subsection*{Progressão e Escalonamento das Ondas}
A progressão da partida fundamenta-se em um modelo dinâmico de escalonamento da dificuldade ao longo das dez ondas planejadas. A cada nova onda completada, incrementam-se de maneira gradual a quantidade total de resíduos gerados, os valores de vida efetiva das unidades e a sua velocidade de avanço pela malha. 

Esse modelo progressivo proporciona uma curva de aprendizado equilibrada, estimulando o jogador a posicionar novas torres e aprimorar as defesas existentes para conter investidas mais densas e velozes. A calibração final dos parâmetros de avanço e geração temporal será conduzida empiricamente durante a integração do ciclo principal de jogo, garantindo estabilidade e fluidez à experiência de jogabilidade.

{\small
\begin{tabularx}{\textwidth}{|p{0.28\textwidth}|X|}
\hline
\textbf{Espécie de Tartaruga} & \textbf{Perfil Tático e Parâmetros Base} \\
\hline
Tartaruga de Couro & Custo 200 conchas; dano 60; alcance 4,0; cadência balanceada. Dispara projéteis de alto impacto direcionados a alvos prioritários de grande resistência. \\
\hline
Tartaruga de Pente & Custo 100 conchas; dano base 20; alcance 2,5; cadência rápida. Implementa mecânica de acerto crítico, duplicando o dano causado a cada quinto ataque consecutivo. \\
\hline
Tartaruga Verde & Custo planejado de 50 conchas; dano 10; alcance 2,0. Unidade defensiva básica de baixo custo, voltada para a contenção inicial nas primeiras ondas. \\
\hline
\end{tabularx}}

Os tipos de resíduos planejados compreendem a \texttt{GarrafaPet} (resistência intermediária e deslocamento moderado), a \texttt{RedeFantasma} (elevada resistência mecânica e velocidade reduzida) e o \texttt{Microplastico} (baixa resistência, porém velocidade acelerada de travessia). A proporção de cada tipo varia conforme a progressão da partida, demandando adaptação contínua da estratégia defensiva do jogador.

\subsection*{Estados da Partida e Apuração de Resultados}
O ciclo operacional da partida divide-se em quatro fases principais:
\begin{enumerate}
    \item \textbf{Preparação:} intervalo entre ondas em que o jogador analisa o tabuleiro e adquire novas torres de tartaruga com suas conchas disponíveis.
    \item \textbf{Combate:} período em que os resíduos deslocam-se pelo trajeto e as torres realizam a detecção de alvos e o disparo de projéteis.
    \item \textbf{Pausa:} suspensão temporária do avanço temporal, interrompendo movimentações e cálculos balísticos sem alteração do estado da partida.
    \item \textbf{Encerramento:} finalização da sessão de jogo por vitória ou derrota, cálculo da bonificação de pontuação vinculada à vida restante da praia e oferta para inserção do registro no ranking.
\end{enumerate}

\clearpage
\section*{3. Diagrama de Classes -- Subsistema de Combate}
O diagrama a seguir sintetiza a arquitetura do subsistema de combate, evidenciando as relações de herança, abstração e associação entre as classes já implementadas na primeira etapa e as extensões planejadas para a conclusão do jogo. Construtores e métodos auxiliares foram omitidos para favorecer a clareza visual.

\begin{center}
\resizebox{0.92\textwidth}{!}{%
\begin{tikzpicture}[x=0.98pt,y=-0.98pt,font=\scriptsize]
\path[use as bounding box] (0,0) rectangle (500,550);
\draw[-{Triangle[open,length=8pt,width=7pt]}] (77,150) -- (77,125) -- (210,125) -- (210,105);
\draw[-{Triangle[open,length=8pt,width=7pt]}] (422,150) -- (422,125) -- (285,125) -- (285,105);
\draw[-{Triangle[open,length=8pt,width=7pt]}] (345,45) -- (325,45);
\draw[-{Triangle[open,length=8pt,width=7pt]}] (155,332) -- (180,332) -- (180,280) -- (77,280) -- (77,260);
\draw[-{Triangle[open,length=8pt,width=7pt]}] (345,332) -- (320,332) -- (320,280) -- (422,280) -- (422,260);
\draw[-{Triangle[open,length=8pt,width=7pt]}] (155,422) -- (180,422) -- (180,280) -- (77,280) -- (77,260);
\draw[-{Triangle[open,length=8pt,width=7pt]}] (345,422) -- (320,422) -- (320,280) -- (422,280) -- (422,260);
\draw[-{Triangle[open,length=8pt,width=7pt]}] (155,512) -- (180,512) -- (180,280) -- (77,280) -- (77,260);
\draw[-{Triangle[open,length=8pt,width=7pt]}] (345,512) -- (320,512) -- (320,280) -- (422,280) -- (422,260);

\draw[fill=gray!5] (170,0) rectangle (325,105);
\node[anchor=north,inner sep=0pt] at (247.5,4) {\textbf{Entidade (abstrata)}};
\draw (170,21) -- (325,21);
\node[anchor=north west,inner sep=0pt] at (175,24) {- posX, posY: int};
\node[anchor=north west,inner sep=0pt] at (175,35) {- ativa: boolean};
\draw (170,51) -- (325,51);
\node[anchor=north west,inner sep=0pt] at (175,57) {+ atualizar()};
\node[anchor=north west,inner sep=0pt] at (175,68) {+ desativar()};

\draw[fill=gray!5] (345,0) rectangle (500,105);
\node[anchor=north,inner sep=0pt] at (422.5,4) {\textbf{Projetil}};
\draw (345,21) -- (500,21);
\node[anchor=north west,inner sep=0pt] at (350,24) {- velocidade: double};
\node[anchor=north west,inner sep=0pt] at (350,35) {- dano: int; tipo: String};
\node[anchor=north west,inner sep=0pt] at (350,46) {- alvo: InimigoLixo};
\draw (345,62) -- (500,62);
\node[anchor=north west,inner sep=0pt] at (350,68) {+ atualizar()};

\draw[fill=gray!5] (0,150) rectangle (155,260);
\node[anchor=north,inner sep=0pt] at (77.5,154) {\textbf{TorreTartaruga (abstrata)}};
\draw (0,171) -- (155,171);
\node[anchor=north west,inner sep=0pt] at (5,174) {- custo, dano: int};
\node[anchor=north west,inner sep=0pt] at (5,185) {- alcance: double};
\node[anchor=north west,inner sep=0pt] at (5,196) {- projetisAtivos: List\textless{}Projetil\textgreater{}};
\draw (0,212) -- (155,212);
\node[anchor=north west,inner sep=0pt] at (5,218) {+ atacar(InimigoLixo)};
\node[anchor=north west,inner sep=0pt] at (5,229) {+ atualizar()};

\draw[fill=gray!5] (345,150) rectangle (500,260);
\node[anchor=north,inner sep=0pt] at (422.5,154) {\textbf{InimigoLixo (abstrata)}};
\draw (345,171) -- (500,171);
\node[anchor=north west,inner sep=0pt] at (350,174) {- vidaMaxima, vidaAtual: int};
\node[anchor=north west,inner sep=0pt] at (350,185) {- velocidade: double};
\node[anchor=north west,inner sep=0pt] at (350,196) {- recompensa: int};
\draw (345,212) -- (500,212);
\node[anchor=north west,inner sep=0pt] at (350,218) {+ receberDano(int)};
\node[anchor=north west,inner sep=0pt] at (350,229) {+ estaDestruido(): boolean};

\draw[fill=gray!5] (0,300) rectangle (155,365);
\node[anchor=north,inner sep=0pt] at (77.5,304) {\textbf{TartarugaVerde (planejada)}};
\draw (0,321) -- (155,321);
\draw (0,329) -- (155,329);
\node[anchor=north west,inner sep=0pt] at (5,335) {+ atacar(InimigoLixo)};

\draw[fill=gray!5] (0,390) rectangle (155,455);
\node[anchor=north,inner sep=0pt] at (77.5,394) {\textbf{TartarugaDeCouro}};
\draw (0,411) -- (155,411);
\draw (0,419) -- (155,419);
\node[anchor=north west,inner sep=0pt] at (5,425) {+ atacar(InimigoLixo)};

\draw[fill=gray!5] (0,480) rectangle (155,545);
\node[anchor=north,inner sep=0pt] at (77.5,484) {\textbf{TartarugaDePente}};
\draw (0,501) -- (155,501);
\node[anchor=north west,inner sep=0pt] at (5,504) {- contadorAtaques: int};
\draw (0,520) -- (155,520);
\node[anchor=north west,inner sep=0pt] at (5,526) {+ atacar(InimigoLixo)};

\draw[fill=gray!5] (345,300) rectangle (500,365);
\node[anchor=north,inner sep=0pt] at (422.5,304) {\textbf{GarrafaPet (planejada)}};
\draw (345,321) -- (500,321);
\draw (345,329) -- (500,329);
\node[anchor=north west,inner sep=0pt] at (350,335) {+ atualizar()};

\draw[fill=gray!5] (345,390) rectangle (500,455);
\node[anchor=north,inner sep=0pt] at (422.5,394) {\textbf{RedeFantasma (planejada)}};
\draw (345,411) -- (500,411);
\draw (345,419) -- (500,419);
\node[anchor=north west,inner sep=0pt] at (350,425) {+ atualizar()};

\draw[fill=gray!5] (345,480) rectangle (500,545);
\node[anchor=north,inner sep=0pt] at (422.5,484) {\textbf{Microplastico (planejado)}};
\draw (345,501) -- (500,501);
\draw (345,509) -- (500,509);
\end{tikzpicture}}
\end{center}

O triângulo com preenchimento em branco indica relação de herança (generalização/especialização). O sinal $-$ representa atributo privado e $+$ indica método público. A classe \texttt{TorreTartaruga} mantém uma coleção de \texttt{Projetil}; cada projétil referencia uma instância de \texttt{InimigoLixo} como alvo de sua trajetória. Os atributos de vida pertencem à hierarquia de \texttt{InimigoLixo}, enquanto a classe abstrata \texttt{Entidade} encapsula posicionamento bidimensional e estado de atividade no mapa.

\clearpage
\section*{4. Diagrama de Classes -- Ciclo de Jogo e Interface}
A coordenação da execução e a interação com o usuário são organizadas de modo a desacoplar a camada de apresentação da lógica de regras do sistema. Assim, a interface com o usuário será concretizada com uma interface gráfica simplificada. O componente \texttt{GameController} atua como mediador entre as ações do usuário e o modelo de domínio, coordenando o \texttt{GameLoop}, o despacho de ondas (\texttt{WaveService}) e as transações de compra de unidades (\texttt{LojaService}).

\begin{center}\begin{tikzpicture}[x=0.98pt,y=-0.98pt,font=\scriptsize]
\path[use as bounding box] (0,0) rectangle (500,520);

\draw[dashed,->] (125,75) -- (125,95) -- (220,95) -- (220,115);
\draw[dashed,->] (375,75) -- (375,95) -- (280,95) -- (280,115);

\draw[dashed,->] (200,215) -- (200,235) -- (77,235) -- (77,260);
\draw[dashed,->] (249,215) -- (249,260);
\draw[-] (327,300) -- (345,300);
\draw[-] (77,355) -- (77,420);
\node[anchor=south west,inner sep=0pt] at (82,395) {1};
\draw[-] (155,315) -- (164,315) -- (164,400) -- (230,400) -- (230,420);
\draw[-] (249,355) -- (249,420);
\node[anchor=south west,inner sep=0pt] at (255,395) {1};
\draw[-] (200,355) -- (200,380) -- (105,380) -- (105,420);

\draw[fill=gray!5] (30,0) rectangle (220,75);
\node[anchor=north,inner sep=0pt] at (125,4) {\textbf{TelaMenuView / Ranking (GUI)}};
\draw (30,21) -- (220,21);
\draw (30,29) -- (220,29);
\node[anchor=north west,inner sep=0pt] at (35,35) {+ exibirMenu()};
\node[anchor=north west,inner sep=0pt] at (35,46) {+ exibirRanking()};
\node[anchor=north west,inner sep=0pt] at (35,57) {+ lerOpcaoUsuario()};

\draw[fill=gray!5] (280,0) rectangle (470,75);
\node[anchor=north,inner sep=0pt] at (375,4) {\textbf{TelaJogoView (GUI)}};
\draw (280,21) -- (470,21);
\draw (280,29) -- (470,29);
\node[anchor=north west,inner sep=0pt] at (285,35) {+ renderizarEstado()};
\node[anchor=north west,inner sep=0pt] at (285,46) {+ capturarAcaoJogador()};
\node[anchor=north west,inner sep=0pt] at (285,57) {+ notificarMensagem(String)};

\draw[fill=gray!5] (172,115) rectangle (327,215);
\node[anchor=north,inner sep=0pt] at (249.5,119) {\textbf{GameController}};
\draw (172,136) -- (327,136);
\draw (172,144) -- (327,144);
\node[anchor=north west,inner sep=0pt] at (177,150) {+ novaPartida()};
\node[anchor=north west,inner sep=0pt] at (177,161) {+ comprarTorre(tipo, col, lin)};
\node[anchor=north west,inner sep=0pt] at (177,172) {+ processarCiclo()};
\node[anchor=north west,inner sep=0pt] at (177,183) {+ pausarPartida()};
\node[anchor=north west,inner sep=0pt] at (177,194) {+ finalizarPartida()};

\draw[fill=gray!5] (0,260) rectangle (155,355);
\node[anchor=north,inner sep=0pt] at (77.5,264) {\textbf{LojaService}};
\draw (0,281) -- (155,281);
\node[anchor=north west,inner sep=0pt] at (5,284) {- jogador: Jogador};
\node[anchor=north west,inner sep=0pt] at (5,295) {- mapa: Mapa};
\draw (0,311) -- (155,311);
\node[anchor=north west,inner sep=0pt] at (5,317) {+ comprar(tipo, col, lin)};
\node[anchor=north west,inner sep=0pt] at (5,328) {+ validarSaldo(custo)};

\draw[fill=gray!5] (172,260) rectangle (327,355);
\node[anchor=north,inner sep=0pt] at (249.5,264) {\textbf{GameLoop}};
\draw (172,281) -- (327,281);
\node[anchor=north west,inner sep=0pt] at (177,284) {- emExecucao: boolean};
\node[anchor=north west,inner sep=0pt] at (177,295) {- inimigos: List\textless{}InimigoLixo\textgreater{}};
\draw (172,311) -- (327,311);
\node[anchor=north west,inner sep=0pt] at (177,317) {+ atualizar()};
\node[anchor=north west,inner sep=0pt] at (177,328) {+ pausar(); retomar()};

\draw[fill=gray!5] (345,260) rectangle (500,355);
\node[anchor=north,inner sep=0pt] at (422.5,264) {\textbf{WaveService}};
\draw (345,281) -- (500,281);
\node[anchor=north west,inner sep=0pt] at (350,284) {- numeroOnda: int};
\draw (345,300) -- (500,300);
\node[anchor=north west,inner sep=0pt] at (350,306) {+ iniciarOnda(int)};
\node[anchor=north west,inner sep=0pt] at (350,317) {+ proximoInimigo()};

\draw[fill=gray!5] (0,420) rectangle (155,515);
\node[anchor=north,inner sep=0pt] at (77.5,424) {\textbf{Jogador}};
\draw (0,441) -- (155,441);
\node[anchor=north west,inner sep=0pt] at (5,444) {- nome: String};
\node[anchor=north west,inner sep=0pt] at (5,455) {- moedas, vida, pontos: int};
\draw (0,471) -- (155,471);
\node[anchor=north west,inner sep=0pt] at (5,477) {+ gastar(int); ganhar(int)};
\node[anchor=north west,inner sep=0pt] at (5,488) {+ perderVida(int)};

\draw[fill=gray!5] (172,420) rectangle (327,515);
\node[anchor=north,inner sep=0pt] at (249.5,424) {\textbf{Mapa}};
\draw (172,441) -- (327,441);
\node[anchor=north west,inner sep=0pt] at (177,444) {- caminho: List\textless{}Coordenada\textgreater{}};
\node[anchor=north west,inner sep=0pt] at (177,455) {- torres: List\textless{}TorreTartaruga\textgreater{}};
\draw (172,471) -- (327,471);
\node[anchor=north west,inner sep=0pt] at (177,477) {+ podeConstruir(col, lin)};
\node[anchor=north west,inner sep=0pt] at (177,488) {+ adicionarTorre(torre)};

\draw[fill=gray!5] (345,420) rectangle (500,515);
\node[anchor=north,inner sep=0pt] at (422.5,424) {\textbf{Exceções de Domínio}};
\draw (345,441) -- (500,441);
\draw (345,449) -- (500,449);
\node[anchor=north west,inner sep=0pt] at (350,455) {PosicaoInvalidaException};
\node[anchor=north west,inner sep=0pt] at (350,466) {RecursosInsuficientesException};
\node[anchor=north west,inner sep=0pt] at (350,477) {(estendem Exception)};
\end{tikzpicture}\end{center}

Seta tracejada indica relação de uso ou dependência funcional; linha contínua expressa associação estrutural. A camada de apresentação é representada por uma interface gráfica simples (JavaFX ou Swing), eliminando interações por terminal e proporcionando navegação intuitiva por telas e botões funcionais. O \texttt{GameLoop} gerencia o estado da partida associando-se a \texttt{Jogador}, \texttt{Mapa} e \texttt{WaveService}. O serviço \texttt{LojaService} valida posições e saldo antes de efetivar modificações nas instâncias do modelo.

\clearpage
\section*{5. Diagrama de Classes -- Persistência de Dados}
O modelo de persistência foi planejado com foco na simplicidade estrutural, empregando arquivos de texto plano locais para armazenamento e dispensando gerenciadores de banco de dados externos. As classes \texttt{GerenciadorRanking} e \texttt{RegistroRanking} já se encontram implementadas e operantes em memória, fazendo uso da interface \texttt{Comparable} para a ordenação automática das pontuações registradas.

Para as etapas subsequentes, projeta-se a classe \texttt{GerenciadorPersistencia}, responsável pelas operações de leitura e escrita dos recordes de pontuação e dos dados essenciais de uma partida em andamento (\texttt{DadosPartida}), registrando informações como a onda alcançada, recursos acumulados e coordenadas das torres ativas.

\begin{center}\begin{tikzpicture}[x=0.98pt,y=-0.98pt,font=\scriptsize]
\path[use as bounding box] (0,0) rectangle (500,340);

\draw[-] (220,75) -- (280,75);
\node[anchor=south west,inner sep=0pt] at (225,68) {1};
\node[anchor=south east,inner sep=0pt] at (275,68) {0..*};

\draw[dashed,->] (120,200) -- (120,130);
\draw[dashed,->] (220,255) -- (280,255);

\draw[fill=gray!5] (20,20) rectangle (220,130);
\node[anchor=north,inner sep=0pt] at (120,24) {\textbf{GerenciadorRanking}};
\draw (20,41) -- (220,41);
\node[anchor=north west,inner sep=0pt] at (25,45) {- records: List\textless{}RegistroRanking\textgreater{}};
\draw (20,68) -- (220,68);
\node[anchor=north west,inner sep=0pt] at (25,74) {+ adicionarRegistro(nome, p, o)};
\node[anchor=north west,inner sep=0pt] at (25,86) {+ getTopRegistros(limite: int): List};

\draw[fill=gray!5] (280,20) rectangle (480,130);
\node[anchor=north,inner sep=0pt] at (380,24) {\textbf{RegistroRanking}};
\draw (280,41) -- (480,41);
\node[anchor=north west,inner sep=0pt] at (285,45) {- nomeJogador: String};
\node[anchor=north west,inner sep=0pt] at (285,57) {- pontuacaoTotal: int};
\node[anchor=north west,inner sep=0pt] at (285,69) {- ondasSobrevividas: int};
\draw (280,88) -- (480,88);
\node[anchor=north west,inner sep=0pt] at (285,94) {+ compareTo(RegistroRanking): int};
\node[anchor=north west,inner sep=0pt] at (285,106) {+ getPontuacaoTotal(): int};

\draw[fill=gray!5] (20,200) rectangle (220,310);
\node[anchor=north,inner sep=0pt] at (120,204) {\textbf{GerenciadorPersistencia (planejada)}};
\draw (20,221) -- (220,221);
\node[anchor=north west,inner sep=0pt] at (25,225) {- caminhoRanking: String};
\node[anchor=north west,inner sep=0pt] at (25,237) {- caminhoPartida: String};
\draw (20,256) -- (220,256);
\node[anchor=north west,inner sep=0pt] at (25,262) {+ salvarRanking(List\textless{}RegistroRanking\textgreater{})};
\node[anchor=north west,inner sep=0pt] at (25,274) {+ carregarRanking(): List};
\node[anchor=north west,inner sep=0pt] at (25,286) {+ salvarPartida(DadosPartida)};
\node[anchor=north west,inner sep=0pt] at (25,298) {+ carregarPartida(): DadosPartida};

\draw[fill=gray!5] (280,200) rectangle (480,310);
\node[anchor=north,inner sep=0pt] at (380,204) {\textbf{DadosPartida (planejada)}};
\draw (280,221) -- (480,221);
\node[anchor=north west,inner sep=0pt] at (285,225) {- nomeJogador: String};
\node[anchor=north west,inner sep=0pt] at (285,237) {- ondaAtual, moedas, vida: int};
\node[anchor=north west,inner sep=0pt] at (285,249) {- torresAtivas: List\textless{}String\textgreater{}};
\draw (280,268) -- (480,268);
\node[anchor=north west,inner sep=0pt] at (285,274) {+ getOndaAtual(): int};
\node[anchor=north west,inner sep=0pt] at (285,286) {+ getTorresAtivas(): List};
\end{tikzpicture}\end{center}

A interface \texttt{Comparable} implementada por \texttt{RegistroRanking} define a ordem natural dos elementos pela pontuação de maneira decrescente. O componente \texttt{GerenciadorPersistencia} utilizará manipuladores de fluxo de entrada e saída padrão do Java (\texttt{BufferedReader} e \texttt{BufferedWriter}) para serializar os registros em arquivos de texto delimitados. O armazenamento é estruturado para salvar e reconstruir posições de torres e recursos sem dependências externas.

\clearpage
\section*{6. Princípios de POO e Organização Modular}
\subsection*{Fluxo de Interação do Usuário}
O fluxo operacional do jogo é estruturado de maneira simples e direta, cobrindo três momentos fundamentais de interação:
\begin{enumerate}
    \item \textbf{Menu Principal:} ponto de acesso inicial onde o jogador pode iniciar uma nova partida, consultar a tabela de melhores pontuações registradas ou encerrar a aplicação.
    \item \textbf{Tela de Jogo:} exibição do mapa e do caminho dos resíduos, acompanhamento contínuo dos indicadores vitais (integridade da praia, saldo de conchas, pontuação e onda vigente), seleção e posicionamento de torres de tartaruga em células válidas, além de comandos para disparo de onda e pausa.
    \item \textbf{Tela de Encerramento:} exibição do resultado de vitória ou derrota, consolidação dos pontos obtidos e formulário para inserção do nome do jogador a fim de salvar o registro no ranking local.
\end{enumerate}

\subsection*{Aplicação dos Pilares da Orientação a Objetos}
A arquitetura do software apoia-se rigorosamente nos conceitos fundamentais da Programação Orientada a Objetos:
\begin{itemize}
    \item \textbf{Abstração (Classes Abstratas e Interfaces):} Concretizada nas classes base \texttt{Entidade}, \texttt{TorreTartaruga} e \texttt{InimigoLixo}, que definem contratos fundamentais de estado e comportamento balístico sem acoplar a implementações finais, bem como no uso da interface nativa \texttt{Comparable} em \texttt{RegistroRanking}.
    \item \textbf{Herança (\textit{is-a}):} Utilização de especialização hierárquica através de \texttt{extends} para as diferentes espécies de defensores (\texttt{TartarugaDeCouro}, \texttt{TartarugaDePente} e \texttt{TartarugaVerde}) e dejetos (\texttt{GarrafaPet}, \texttt{RedeFantasma}), compartilhando atributos vitais e de posicionamento.
    \item \textbf{Composição (\textit{has-a}):} Emprega-se a composição estrutural favorecendo o baixo acoplamento em diversos pontos críticos do sistema:
    \begin{itemize}
        \item \texttt{TorreTartaruga} \textit{tem-uma} coleção de projéteis ativos (\texttt{List<Projetil>});
        \item \texttt{GerenciadorRanking} \textit{tem-uma} lista de recordes de pontuação (\texttt{List<RegistroRanking>});
        \item \texttt{Projetil} \textit{tem-um} alvo do tipo \texttt{InimigoLixo} que persegue até o impacto.
    \end{itemize}
    \item \textbf{Polimorfismo e Ligação Tardia:} Manifestado nas chamadas ao método abstrato \texttt{atacar(InimigoLixo)}, onde cada espécie concreta de tartaruga executa dinamicamente seu algoritmo de disparo (dano concentrado versus contagem para acerto crítico), sem que o despachante precise averiguar a subclasse em tempo de compilação.
    \item \textbf{Sobrescrita (\textit{Override}) versus Sobrecarga (\textit{Overload}):}
    \begin{itemize}
        \item \textit{Sobrescrita (\texttt{@Override}):} Modificação de comportamento de métodos herdados, como \texttt{atualizar()} em \texttt{TorreTartaruga} e \texttt{Projetil}, e \texttt{compareTo(RegistroRanking)} para inversão da ordenação natural;
        \item \textit{Sobrecarga (\textit{Overload}):} Construtores com assinaturas e listas de parâmetros distintas em \texttt{Entidade} (\texttt{Entidade()} padrão versus \texttt{Entidade(posX, posY)} parametrizado), permitindo instanciação flexível no catálogo de compras antes do posicionamento em malha.
    \end{itemize}
    \item \textbf{Encapsulamento e Modularidade:} Todos os atributos de integridade, custo, dano e coordenadas foram declarados sob modificadores de acesso restrito (\texttt{private}), com acesso exclusivamente controlado por \textit{getters} e \textit{setters} munidos de validação.
    \item \textbf{Tratamento de Exceções:} Classes de exceção de domínio personalizadas (\texttt{PosicaoInvalidaException} e \texttt{RecursosInsuficientesException}) que impedem jogadas ilegais ou inconsistências de saldo através de blocos \texttt{try/catch}.
\end{itemize}

\subsection*{Estrutura de Pacotes da Aplicação}
Em conformidade com a convenção modular exigida na disciplina, o código-fonte foi distribuído nos seguintes pacotes:
\begin{itemize}
    \item \texttt{br.ufv.ccf313.projeto}: Contém a classe executável \texttt{Main};
    \item \texttt{br.ufv.ccf313.projeto.model}: Entidades e modelos de domínio (\texttt{Entidade}, \texttt{TorreTartaruga}, \texttt{TartarugaDeCouro}, \texttt{TartarugaDePente}, \texttt{InimigoLixo}, \texttt{Projetil}, \texttt{RegistroRanking});
    \item \texttt{br.ufv.ccf313.projeto.service}: Regras de negócio e serviços (\texttt{GerenciadorRanking}, \texttt{GameLoopService}, \texttt{LojaService}, \texttt{WaveService});
    \item \texttt{br.ufv.ccf313.projeto.repository}: Camada de persistência física em arquivos de texto plano;
    \item \texttt{br.ufv.ccf313.projeto.controller}: Mediação de comandos entre a interface e os serviços de domínio;
    \item \texttt{br.ufv.ccf313.projeto.view}: Classes da interface gráfica do usuário (planejada em JavaFX / Swing);
    \item \texttt{br.ufv.ccf313.projeto.test}: Baterias de testes unitários automatizados em JUnit.
\end{itemize}

\subsection*{Divisão de Responsabilidades da Equipe}
A Tabela~\ref{tab:responsaveis} detalha a distribuição das atribuições entre os integrantes do grupo, garantindo contribuições colaborativas em módulos específicos:

\begin{table}[h!]
\centering
\small
\begin{tabularx}{\textwidth}{|p{0.24\textwidth}|p{0.25\textwidth}|X|}
\hline
\textbf{Integrante} & \textbf{Módulo de Atuação} & \textbf{Classes / Responsabilidades Principais} \\
\hline
1. Lethícia Gonçalves da Silva & Base do Combate & Cria a classe abstrata \texttt{Entidade.java} e a lógica matemática de posições e colisões no mapa (\texttt{CollisionService.java}). \\
\hline
2. Tomás Viana Resende & Tartarugas / Torres & Cria a classe abstrata \texttt{TorreTartaruga.java}, tipos de tartarugas (\texttt{TartarugaDeCouro}, \texttt{TartarugaDePente}, \texttt{TartarugaVerde}), projéteis e mecânica de upgrades. \\
\hline
3. Bruna Lima Ribeiro & Lixos / Inimigos & Cria a classe abstrata \texttt{InimigoLixo.java} e os tipos de lixo marinho (\texttt{GarrafaPet}, \texttt{RedeFantasma}, \texttt{Microplastico}) com escalonamento de vida por round. \\
\hline
4. Helena Santos Perez & Motor do Jogo \& Exceções & Cria o loop principal da partida (\texttt{GameLoop.java}), controle das ondas de inimigos (\texttt{WaveService.java}) e as exceções customizadas (\texttt{PosicaoInvalidaException}, \texttt{RecursosInsuficientesException}). \\
\hline
5. Otávio Soares Pedroso & Mapa \& Jogador & Cria a lógica do mapa/caminho do lixo (\texttt{Mapa.java}, \texttt{CelulaGrid.java}), gestão de moedas, vida da praia (\texttt{Jogador.java}) e as regras de compra da loja (\texttt{LojaService.java}). \\
\hline
6. Pedro Mendes Gonzaga & Persistência \& Arquivos & Cria as classes de persistência para salvar e carregar dados em arquivos de texto (saves de partidas, ranking e configurações). \\
\hline
7. Marina Larissa Xavier Soares & Interface da Partida \& Controller & Cria a tela gráfica principal do jogo (JavaFX/Swing) com a renderização do mapa, animações e cliques do mouse (\texttt{GameController.java}, \texttt{TelaJogoView.java}). \\
\hline
8. Caroline Melo Silva & Menus \& Testes Unitários & Cria as telas de menu/ranking do sistema (\texttt{TelaMenuView.java}, \texttt{TelaRankingView.java}) e escreve a suíte de testes com JUnit para validação das regras. \\
\hline
\end{tabularx}
\caption{Quadro de divisão de responsabilidades da equipe.}
\label{tab:responsaveis}
\end{table}

\clearpage
\section*{7. Planejamento das 5 Funcionalidades Completas (CRUD)}
Conforme as diretrizes da disciplina, o sistema contempla o planejamento de cinco funcionalidades completas, cada uma provendo operações de criação (Create), consulta (Read), atualização (Update) e remoção (Delete):

\begin{enumerate}
    \item \textbf{Gerenciamento de Ranking e Recordes (Implementada na Entrega Parcial):}
    \begin{itemize}
        \item \textit{Create:} Inserção de novos registros de pontuação ao concluir partidas (\texttt{adicionarRegistro}).
        \item \textit{Read:} Consulta e ordenação das maiores pontuações (\texttt{getTopRegistros}) e busca individual por nome (\texttt{buscarRegistro}).
        \item \textit{Update:} Atualização da pontuação e ondas sobrevividas de um jogador (\texttt{atualizarRegistro}).
        \item \textit{Delete:} Exclusão de registros de pontuação obsoletos (\texttt{removerRegistro}).
    \end{itemize}

    \item \textbf{Gerenciamento de Torres no Tabuleiro (Planejada):}
    \begin{itemize}
        \item \textit{Create:} Aquisição e posicionamento de novas torres de tartaruga em células válidas do grid.
        \item \textit{Read:} Inspeção visual e textual das propriedades, dano e raio de alcance de cada torre ativa.
        \item \textit{Update:} Evolução de nível (\textit{upgrade}) da torre, aprimorando atributos de cadência e dano.
        \item \textit{Delete:} Venda e remoção da torre do tabuleiro com reembolso proporcional em conchas.
    \end{itemize}

    \item \textbf{Gerenciamento de Jogador e Perfis (Planejada):}
    \begin{itemize}
        \item \textit{Create:} Cadastro de perfil de jogador contendo nome e identificador.
        \item \textit{Read:} Consulta de dados cadastrais, saldo de conchas acumulado e estatísticas de vitórias.
        \item \textit{Update:} Edição de dados do perfil e preferências de exibição.
        \item \textit{Delete:} Exclusão de contas de perfil e histórico associado.
    \end{itemize}

    \item \textbf{Gerenciamento de Ondas de Resíduos (Planejada):}
    \begin{itemize}
        \item \textit{Create:} Criação e configuração de novas ondas com listas específicas de dejetos marinhos.
        \item \textit{Read:} Consulta da relação de resíduos e temporização de chegada da onda em vigor.
        \item \textit{Update:} Ajuste de dificuldade e parâmetros de escalonamento dos dejetos.
        \item \textit{Delete:} Cancelamento ou redefinição da fila de ondas configuradas.
    \end{itemize}

    \item \textbf{Gerenciamento de Partidas Salvas (Save Game) (Planejada):}
    \begin{itemize}
        \item \textit{Create:} Gravação do estado corrente da partida em slot de arquivo texto (\texttt{DadosPartida}).
        \item \textit{Read:} Listagem de partidas salvas e leitura para retomada do jogo.
        \item \textit{Update:} Sobrescrita do progresso da sessão com dados atualizados da onda vigente.
        \item \textit{Delete:} Exclusão de arquivos de salvamento antigos.
    \end{itemize}
\end{enumerate}

\section*{8. Resultados da Entrega Parcial e Relatório de Testes}
\subsection*{Classe Principal Executável (\texttt{Main})}
Para viabilizar a verificação imediata dos artefatos da primeira entrega sem a dependência da interface gráfica final, implementou-se a classe executável \texttt{Main} no pacote principal. Esta classe atua como ponto de entrada para:
\begin{enumerate}
    \item Demonstrar a instanciação polimórfica das classes de combate herdadas de \texttt{Entidade} e \texttt{TorreTartaruga};
    \item Executar em console o ciclo integral das quatro operações fundamentais do CRUD de Ranking (inserção, consulta, atualização de pontuação e remoção).
\end{enumerate}

\subsection*{Relatório Textual dos Testes Automatizados (JUnit)}
Em conformidade com a seção de testes da disciplina, elaborou-se a suíte de testes unitários \texttt{GerenciadorRankingTest}, cobrindo a funcionalidade completa implementada:
\begin{itemize}
    \item \textbf{Partes testadas:} Métodos \texttt{adicionarRegistro}, \texttt{getTopRegistros}, \texttt{buscarRegistro}, \texttt{atualizarRegistro} e \texttt{removerRegistro} da classe \texttt{GerenciadorRanking}.
    \item \textbf{Cenários considerados:}
    \begin{enumerate}
        \item Inserção de múltiplos registros desordenados e conferência da ordenação decrescente automática;
        \item Busca por nome existente (retorno correto) e por nome não cadastrado (retorno \texttt{null});
        \item Atualização de pontuação e ondas de jogador cadastrado (sucesso) e tentativa em registro inexistente (\texttt{false});
        \item Remoção de registro existente (sucesso e decréscimo do total) e tentativa em registro inexistente;
        \item Validação de integridade: bloqueio de nomes vazios ou nulos via disparo de \texttt{IllegalArgumentException}.
    \end{enumerate}
    \item \textbf{Resultados esperados:} Sucesso em todas as asserções de igualdade (\texttt{assertEquals}), não-nulos (\texttt{assertNotNull}) e tratamento de exceção (\texttt{assertThrows}).
    \item \textbf{Resultados obtidos:} 100\% de êxito na execução dos cinco casos de teste automatizados sob JUnit 5.
\end{itemize}

\clearpage
\section*{9. Conclusão e Trabalhos Futuros}
\subsection*{Considerações da Primeira Etapa}
A primeira etapa de desenvolvimento do projeto \textit{Tartarugas vs. Lixo Marinho} cumpriu com êxito os requisitos estruturais e de modelagem propostos. Estabeleceu-se uma arquitetura sólida orientada a objetos, comprovada pela implementação das classes basilares de combate (\texttt{Entidade}, \texttt{TorreTartaruga}, \texttt{InimigoLixo}), especializações concretas de defensores (\texttt{TartarugaDeCouro} e \texttt{TartarugaDePente}) e a entidade balística \texttt{Projetil}.

Adicionalmente, consolidou-se o subsistema de pontuação e ranking em memória com suporte a ordenação automática baseada na interface \texttt{Comparable}, validada por suíte de testes unitários automatizados com o \textit{framework} JUnit. A modelagem em diagramas de classes proporcionou uma visão arquitetural abrangente e garantiu o desacoplamento entre os elementos de combate, a camada de regras da partida e as estruturas de persistência e visualização.

\subsection*{Cronograma de Trabalhos Futuros}
O planejamento para a segunda etapa do projeto contempla o desenvolvimento e a integração dos módulos restantes, seguindo o roteiro técnico delineado abaixo:
\begin{enumerate}
    \item \textbf{Conclusão das Unidades de Domínio:} implementação da especialização inicial \texttt{TartarugaVerde} e das classes concretas de dejetos (\texttt{GarrafaPet}, \texttt{RedeFantasma} e \texttt{Microplastico}), calibrando seus atributos de movimentação e resistência.
    \item \textbf{Construção do Ciclo de Jogo e Ondas:} desenvolvimento do motor contínuo (\texttt{GameLoop}) e do despachante \texttt{WaveService}, integrando a geração temporal dos resíduos, a navegação pelo trajeto da malha e a sincronização do avanço de projéteis e verificação de colisões.
    \item \textbf{Regras Econômicas e Tabuleiro:} integração das classes \texttt{Mapa}, \texttt{Jogador} e \texttt{LojaService}, com validação rigorosa de posições e recursos via exceções de domínio (\texttt{PosicaoInvalidaException} e \texttt{RecursosInsuficientesException}).
    \item \textbf{Camada de Interface com o Usuário:} construção da interface gráfica simples (JavaFX / Swing) mediada por \texttt{GameController}, garantindo a renderização do mapa, animações visuais e o controle de ações do jogador via cliques de mouse.
    \item \textbf{Persistência Simplificada em Arquivo:} desenvolvimento da classe \texttt{GerenciadorPersistencia} para gravação e leitura de recordes e recuperação de partidas interrompidas a partir de arquivos de texto locais.
\end{enumerate}

\subsection*{Considerações Finais}
A abordagem modular e os fundamentos de Orientação a Objetos adotados asseguram viabilidade técnica para a conclusão do projeto dentro do cronograma da disciplina. A arquitetura estabelecida combina rigor conceitual com flexibilidade operacional, permitindo que a aplicação final ofereça uma experiência de uso intuitiva e atrativa por meio de interface gráfica, mantendo o objetivo central de sensibilização ambiental por meio da tecnologia.

\end{document}
